% ============================================================================
%  preamble.tex —— 导言区（宏包、版式、定理环境、记号宏）
%  项目：Frenkel–Witten, Geometric Endoscopy and Mirror Symmetry
%  编译：xelatex main.tex （连续两遍）
% ============================================================================

% !TEX program = xelatex
\usepackage[UTF8,scheme=plain,fontset=fandol]{ctex}

% ---------------------------------------------------------------------------
% 数学与排版
% ---------------------------------------------------------------------------
\usepackage{amsmath}
\usepackage{amssymb}
\usepackage{amsthm}
\usepackage{mathtools}
\usepackage{bm}
% 注意：不载入 amscd，因为宏 \CD 用作 D-模记号；交换图请用 array 环境重建（见 SPEC §4）。
\usepackage[a4paper,top=2.4cm,bottom=2.4cm,left=2.3cm,right=2.3cm]{geometry}
\usepackage{graphicx}
\usepackage{xcolor}
\usepackage{enumitem}
\usepackage{booktabs}
\usepackage{array}
\usepackage{caption}
\usepackage{float}
\usepackage{fancyhdr}
\usepackage{tikz}
\usetikzlibrary{arrows.meta}
\usepackage[hidelinks,bookmarksnumbered=true,bookmarksopen=true]{hyperref}

\hypersetup{
    hypertexnames=false,   % 显式 \tag 编号与自动计数器易产生重名锚点，故关闭
    pdftitle={几何内窥与镜像对称（中文译注本）},
    pdfauthor={Edward Frenkel, Edward Witten（原著）；中文译注},
    pdfsubject={arXiv:0710.5939v3 [math.AG]},
    pdfkeywords={geometric Langlands, endoscopy, mirror symmetry, Hitchin fibration,
                 A-brane, D-module, Hecke eigensheaf, L-packet}
}

% ---------------------------------------------------------------------------
% 版式
% ---------------------------------------------------------------------------
\allowdisplaybreaks
\renewcommand{\contentsname}{目\quad 录}
\renewcommand{\abstractname}{摘\quad 要}
\renewcommand{\refname}{参考文献}
\renewcommand{\figurename}{图}
\renewcommand{\tablename}{表}
\setlength{\emergencystretch}{2em}
\setlength{\parskip}{0.2em}
\linespread{1.06}

\makeatletter
\renewcommand*\l@section{\@dottedtocline{1}{0em}{1.6em}}
\makeatother

% 定理类环境：编号一律照原文手写，避免自动计数与原文不一致
\newenvironment{namedthm}[1]
  {\par\medskip\noindent\textbf{定理 #1.}\itshape\ignorespaces}
  {\par\medskip}
\newenvironment{namedprop}[1]
  {\par\medskip\noindent\textbf{命题 #1.}\itshape\ignorespaces}
  {\par\medskip}
\newenvironment{namedlemma}[1]
  {\par\medskip\noindent\textbf{引理 #1.}\itshape\ignorespaces}
  {\par\medskip}
\newenvironment{namedcor}[1]
  {\par\medskip\noindent\textbf{推论 #1.}\itshape\ignorespaces}
  {\par\medskip}
\newenvironment{namedconj}[1]
  {\par\medskip\noindent\textbf{猜想 #1.}\itshape\ignorespaces}
  {\par\medskip}
\newenvironment{nameddef}[1]
  {\par\medskip\noindent\textbf{定义 #1.}\upshape\ignorespaces}
  {\par\medskip}

% 「本节重点」区块
\newenvironment{keytakeaway}
  {\par\medskip\begin{quote}\small\noindent
   \rule{\linewidth}{0.6pt}\par\nopagebreak\smallskip\noindent
   \textbf{本节重点}\par\nobreak\smallskip\noindent\ignorespaces}
  {\par\smallskip\noindent\rule{\linewidth}{0.6pt}\end{quote}\medskip}

% 「译者说明」区块
\newenvironment{translatornote}
  {\begin{quote}\small\noindent\ignorespaces}
  {\end{quote}\medskip}

% 页眉页脚
\pagestyle{fancy}
\fancyhf{}
\fancyhead[L]{\small 几何内窥与镜像对称}
\fancyhead[R]{\small Frenkel--Witten（中文译注本）}
\fancyfoot[C]{\thepage}
\renewcommand{\headrulewidth}{0.4pt}

% 记号宏（与 GLOSSARY 一致；不要新增）
\newcommand{\Bun}{\mathrm{Bun}}
\newcommand{\CD}{\mathcal{D}}
\newcommand{\CM}{\mathcal{M}}
\newcommand{\CH}{\mathcal{H}}
\newcommand{\CC}{\mathbb{C}}
\newcommand{\ZZ}{\mathbb{Z}}
\newcommand{\RR}{\mathbb{R}}
\newcommand{\PP}{\mathbb{P}}
\newcommand{\Tr}{\operatorname{Tr}}
\newcommand{\tr}{\operatorname{tr}}
\newcommand{\dd}{\mathrm{d}}
\newcommand{\ii}{\mathrm{i}}

% 分片容错：尚未交付时不中断编译
\newcommand{\frag}[1]{%
  \IfFileExists{sections/#1.tex}{\input{sections/#1.tex}}{%
    \par\noindent\textcolor{red}{[缺分片：#1]}\par}}
